\subsection{Isolated subgroups of an ordered group}

To study the situation in no. 1 from the point of view of valuations we shall need Definition 1 and Proposition 3 below.

DEFINITION 1. A subgroup H of an ordered group G is called isolated if the relations \(0 \leqslant y \leqslant x\) and \(x \in \mathbf{H}\) imply \(y \in \mathbf{H}\) .

Example (1) Let A and B be two ordered groups; let \(A \times B\) be given the lexicographic order (i.e.~`` \((a, b) \leqslant (a', b')\) '' is equivalent to `` \((a < a')\) or \((a = a' \text{ and } b \leqslant b')\) ''). The second factor B of \(A \times B\) is then, as is seen immediately, an isolated subgroup of \(A \times B\) .

PROPOSITION 3. Let G be an ordered group and P the set of its positive elements.

\begin{enumerate}
\def\labelenumi{(\alph{enumi})}
\item
  The kernel of an increasing homomorphism of G to an ordered group is an isolated subgroup of G.
\item
  Conversely, let H be an isolated subgroup of G and g the canonical homomorphism of G onto G/H. Then \(g(P)\) is the set of positive elements of an ordered group structure on G/H. Moreover, if G is totally ordered, so is G/H.
\item
  Let \(f\) be an increasing homomorphism from \(G\) to an ordered group; let \(H\) denote the kernel of \(f\) . If \(0 \leqslant y \leqslant x\) and \(x \in H\) , then \(0 \leqslant f(y) \leqslant f(x) = 0\) , whence \(f(y) = 0\) , that is \(y \in H\) . Hence \(H\) is isolated.
\item
  Let \(\mathbf{H}\) be an isolated subgroup of \(\mathbf{G}\) and \(g\colon \mathbf{G}\to \mathbf{G} / \mathbf{H}\) . Let \(\mathbf{P}' = g(\mathbf{P})\) . Clearly \(\mathbf{P}' + \mathbf{P}'\subset \mathbf{P}'\) . Also
\end{enumerate}

\[
\mathrm{P} ^ {\prime} \cap (- \mathrm{P} ^ {\prime}) = \{0 \},
\]

for, if \(x\) and \(y\) are elements of \(\mathbf{P}\) such that \(g(x) = -g(y)\) , then \(x + y \in \mathbf{H}\) , whence \(x \in \mathbf{H}\) and \(y \in \mathbf{H}\) since \(\mathbf{H}\) is isolated; hence \(g(x) = g(y) = 0\) . Thus \(\mathbf{P}'\) is the set of positive elements of an ordered group structure on \(\mathbf{G} / \mathbf{H}\) (Algebra, Chapter VI, §1, no. 3, Proposition 3). Finally, if \(\mathbf{G}\) is totally ordered, then \(\mathbf{P} \cap (-\mathbf{P}) = \mathbf{G}\) , whence \(\mathbf{P}' \cup (-\mathbf{P}') = \mathbf{G} / \mathbf{H}\) and therefore \(\mathbf{G} / \mathbf{H}\) is totally ordered (loc. cit.).

Example (2) If we reconsider the example where G is a lexicographic product \(A \times B\) and H = B, the ordered group G/H is canonically identified with A.

